% !PRINTED 472
\weberpage{472}

Die Gleichung (22) hat aber wieder eine Gleichung der Form (21) zur Folge, so daß in beiden die Vorzeichen übereinstimmen, und die Gleichung (21) fordert nach § 40, (7):
\weq{23}{
\begin{aligned}
&\frac{c + \partial\omega}{a} = \frac{\gamma + \delta\omega}{\alpha + \beta\omega}, \quad \delta \equiv 0\ (\mathrm{mod}\ 2), \\
&\gamma^2 - 1 + \gamma(2\alpha - \delta) \equiv 0 \ \text{oder} \equiv 8 \quad (\mathrm{mod}\ 16),
\end{aligned}
}
und zwar gilt, je nachdem die eine oder die andere Kongruenz stattfindet, in (21) und (22) das eine oder das andere Vorzeichen.

Nehmen wir $m = n$, also $m \equiv -1\ (\mathrm{mod}\ 8)$, so ergibt sich aus (2) und (23) ganz wie oben:
\[
\begin{aligned}
\beta &= A, & n\gamma &= -C + cB, \\
\alpha &= B, & n\delta &= -B + cA,
\end{aligned}
\]
also aus der zweiten Kongruenz (23):
\[
C^2 + CB - 1 \equiv 0 \ \text{oder} \equiv 8 \quad (\mathrm{mod}\ 16),
\]
woraus zunächst folgt, daß $B$ jedenfalls durch 8 teilbar sein muß; vermehren wir aber $B$ um ein ungerades Vielfache von 8, so geht die eine dieser Kongruenzen in die andere über, und infolgedessen geht in (21) das eine in das andere Vorzeichen über. Für ein bestimmtes $\omega$ besteht also nur die eine der beiden Gleichungen (22) und es folgt:

7. \gesperrt{Ist $m \equiv 7\ (\mathrm{mod}\ 8)$, so ist $\sqrt{2}\, f(\omega)^3$ Klasseninvariante.}

Wenn wir einer durch alle bekannten Beispiele bestätigten Induktion vertrauen dürfen, so besteht noch das folgende Theorem:

8. \gesperrt{Ist $m \equiv 3\ (\mathrm{mod}\ 8)$, so ist $f(\omega)^3$ Klasseninvariante.}

Indessen fehlt hierfür noch der allgemeine Beweis.

Ist $m \equiv 2\ (\mathrm{mod}\ 4)$, so wenden wir dasselbe Verfahren an, wie oben im Falle $m \equiv 1\ (\mathrm{mod}\ 8)$.

Wir setzen:
\weq{24}{2n = m + y^2,}
und nehmen $y = 0$ oder $= \pm 2$, so daß $n \equiv 1\ (\mathrm{mod}\ 4)$ wird. % [sic] printed second equality omits y

Wenn wir dann in der Gleichung (19)
\weq{25}{v_1 = f_2\left(\frac{\omega}{2}\right)^3 = \frac{\sqrt{2}^{3}}{f_1(\omega)^3}, \quad u_1 = f_1(\omega)^3}
setzen, so ergibt sich eine Gleichung, die nur $\sqrt{2}\, f_1(\omega)^6$ und rationale Koeffizienten enthält. (19) ist aber nur dann erfüllt [§ 40, (7)], wenn
%CONT>>
